% Performance evaluation section for the report.
% Assumptions:
% 1. The main report preamble loads `graphicx`.
% 2. If the report is compiled from a different working directory, adjust the
%    figure paths below accordingly.

\section{Performance Evaluation}
\label{sec:performance}

\subsection{Experimental Set-up}

We evaluate the proposed LLM-based SQL optimization pipeline on the TPC-DS decision-support benchmark using PostgreSQL with the SF=1 dataset. The workload is drawn from the 99 benchmark queries in the testing set. Each query is executed against the populated PostgreSQL instance to support both semantic validation and physical performance measurement. Since the objective is not merely to generate alternative SQL text, but to generate rewrites that preserve semantics while improving execution behavior, the evaluation combines correctness checks with runtime-oriented and resource-oriented metrics.

The evaluation is organized into two phases. In Phase~1, a controlled screening is performed on 10 randomly sampled queries to compare prompt and reasoning combinations. In the collected run, this produced 12 combinations corresponding to four prompt strategies ($P0$--$P3$) and three reasoning strategies ($R0$--$R2$), all instantiated with \texttt{gpt-5-mini}. In Phase~2, the best Phase~1 configuration is applied to a larger sample of 30 TPC-DS queries. In both phases, three candidate rewrites are generated per query and each candidate is benchmarked using five execution trials.

Semantic validity is enforced before any speed claim is accepted. For valid candidates, execution performance is measured using the median execution time reported by \texttt{EXPLAIN (ANALYZE, BUFFERS, FORMAT JSON)}. Candidate speedup is defined as the ratio between baseline and candidate execution time. Resource efficiency is summarized using a derived memory score based on PostgreSQL buffer and temporary I/O behavior. Final ranking uses a composite score with weights $0.7$ for speedup and $0.3$ for memory score. Because the collected results contain several large outliers, medians are treated as the primary summary statistic.

\begin{table}[htbp]
\centering
\caption{Summary of the experimental configuration used in the reported runs.}
\label{tab:exp-setup}
\begin{tabular}{|l|l|}
\textbf{Aspect} & \textbf{Configuration} \\
\hline
Dataset & TPC-DS SF=1 on PostgreSQL \\
Query pool & 99 analytical SQL queries \\
Phase~1 workload & 10 randomly sampled queries \\
Phase~2 workload & 30 randomly sampled queries \\
Candidate count per run & 3 \\
Execution trials per query & 5 \\
Phase~1 search space & $P0$--$P3$ $\times$ $R0$--$R2$ using \texttt{gpt-5-mini} \\
Best configuration selected & \texttt{P3\_\_R0\_\_gpt-5-mini} \\
Correctness metric & semantic equivalence via result-set validation \\
Performance metrics & execution time, planning time, speedup, memory score \\
Ranking score & $0.7 \times$ speedup $+ 0.3 \times$ memory score \\
\hline
\end{tabular}
\end{table}

\subsection{Phase~1: Comparative Screening of Prompt/Reasoning Combinations}

Phase~1 asks which prompt/reasoning combination provides the best balance between correctness and performance. The macro-level comparison is shown in Figure~\ref{fig:phase1-leaderboard}, while Figure~\ref{fig:phase1-heatmap} summarizes the interaction between prompt strategy and reasoning strategy.

At the combination level, the best configuration is \texttt{P3\_\_R0\_\_gpt-5-mini}. It achieves a valid selection rate of 90\% and the highest median selected speedup among all screened combinations, at 1.267$\times$. The next strongest configuration is \texttt{P2\_\_R1\_\_gpt-5-mini}, which maintains the same validity but a slightly lower median speedup of 1.213$\times$. Weaker configurations such as \texttt{P1\_\_R0\_\_gpt-5-mini} and \texttt{P1\_\_R1\_\_gpt-5-mini} remain close to performance parity, suggesting that engine-aware prompting alone is insufficient to drive consistent gains.

Viewed by prompt strategy, $P3$ yields the highest median speedup (1.250$\times$), followed by $P2$ (1.114$\times$). This suggests that richer schema and statistics context helps the model identify more aggressive and more beneficial rewrites. However, the gain comes with a trade-off: $P3$ has a lower overall valid selection rate (76.7\%) than $P2$ (86.7\%). Thus, richer context appears to increase optimization power, but also increases the chance of invalid or unstable rewrites when combined with less suitable reasoning modes.

Viewed by reasoning strategy, $R2$ (two-pass prompting) achieves the highest median speedup (1.117$\times$), but its valid selection rate is only 65\%, substantially below the 90\% achieved by both $R0$ and $R1$. This indicates that explicit multi-step prompting can improve aggressiveness, but also increases the risk of semantically unsafe outputs. In contrast, $R0$ and $R1$ are much more reliable, with $R0$ ultimately forming the best overall combination when paired with $P3$.

Taken together, the Phase~1 results justify the selection of \texttt{P3\_\_R0\_\_gpt-5-mini} for the larger-scale evaluation. It preserves the high validity of direct prompting while benefiting from richer structural context.

\begin{table}[htbp]
\centering
\caption{Top Phase~1 combinations ranked by validity first and median selected speedup second.}
\label{tab:phase1-top-combos}
\begin{tabular}{|l|c|c|c|c|}
\textbf{Combo} & \textbf{Valid rate} & \textbf{Median speedup} & \textbf{Median memory score} & \textbf{Missing speedup rate} \\
\hline
\texttt{P3\_\_R0\_\_gpt-5-mini} & 90.0\% & 1.267 & 0.420 & 10.0\% \\
\texttt{P2\_\_R1\_\_gpt-5-mini} & 90.0\% & 1.213 & 0.417 & 10.0\% \\
\texttt{P0\_\_R0\_\_gpt-5-mini} & 90.0\% & 1.021 & 0.432 & 10.0\% \\
\texttt{P0\_\_R1\_\_gpt-5-mini} & 90.0\% & 1.002 & 0.422 & 0.0\% \\
\texttt{P1\_\_R0\_\_gpt-5-mini} & 90.0\% & 0.996 & 0.425 & 0.0\% \\
\hline
\end{tabular}
\end{table}

\begin{figure}[htbp]
\centering
\includegraphics[width=0.90\linewidth]{datasets/artifacts/findings_figures/phase1_combo_leaderboard.png}
\caption{Phase~1 combination leaderboard. Bars report median selected speedup, while the overlaid line reports valid selection rate.}
\label{fig:phase1-leaderboard}
\end{figure}

\begin{figure}[htbp]
\centering
\includegraphics[width=0.72\linewidth]{datasets/artifacts/findings_figures/phase1_prompt_reasoning_heatmap.png}
\caption{Median selected speedup in Phase~1 as a function of prompt strategy and reasoning strategy.}
\label{fig:phase1-heatmap}
\end{figure}

\subsection{Phase~2: Evaluation of the Best Configuration}

Phase~2 applies the selected configuration, \texttt{P3\_\_R0\_\_gpt-5-mini}, to a larger set of 30 randomly sampled TPC-DS queries. The aggregate picture is mixed. Although the system always produces a ranked selection, only 24 of the 30 selected candidates are semantically valid, yielding a valid selection rate of 80\%. Across the valid candidates, the median speedup is 1.002$\times$, indicating near-parity performance in the typical case rather than broad-based large improvement. The mean speedup is much larger at 5.359$\times$, but this average is inflated by a few extreme wins and therefore is not representative of the central tendency.

In outcome terms, five queries show strong improvement, nine show mild improvement, three are near parity, six regress, one valid query lacks a computable speedup because of missing baseline benchmark timing, and six selected candidates are invalid. Hence, the best Phase~1 configuration generalizes only partially: it produces substantial improvements on some queries, but it does not improve the workload uniformly.

\begin{table}[htbp]
\centering
\caption{Aggregate Phase~2 results for the selected best configuration \texttt{P3\_\_R0\_\_gpt-5-mini}.}
\label{tab:phase2-summary}
\begin{tabular}{|l|c|}
\textbf{Metric} & \textbf{Value} \\
\hline
Queries evaluated & 30 \\
Selected candidates & 30 \\
Semantically valid selected candidates & 24 \\
Valid selection rate & 80.0\% \\
Median speedup & 1.002$\times$ \\
Mean speedup & 5.359$\times$ \\
25th percentile speedup & 0.956$\times$ \\
75th percentile speedup & 1.066$\times$ \\
Median memory score & 0.416 \\
Strong improvements ($\geq 1.10\times$) & 5 \\
Mild improvements ($1.00\times$--$1.10\times$) & 9 \\
Near parity ($0.95\times$--$1.00\times$) & 3 \\
Regressions ($< 0.95\times$) & 6 \\
Missing speedup & 1 \\
Invalid selections & 6 \\
\hline
\end{tabular}
\end{table}

\subsection{Why the Best Configuration Helps or Hurts}

The micro-level results show that the best configuration is highly effective on a small number of queries but not uniformly effective across the workload. The strongest improvement occurs on \texttt{query\_81}, which achieves a speedup of 94.66$\times$, followed by \texttt{query\_31} at 7.47$\times$, \texttt{query\_69} at 1.52$\times$, \texttt{query\_54} at 1.32$\times$, and \texttt{query\_83} at 1.28$\times$. At the other end of the spectrum, the worst regression occurs on \texttt{query\_72} (0.493$\times$), followed by \texttt{query\_51} (0.750$\times$) and \texttt{query\_71} (0.783$\times$). This spread indicates that the chosen prompting strategy can discover highly beneficial rewrites when the original query contains a favorable optimization opportunity, but the same strategy can also over-structure queries that are already close to a good physical form.

Figure~\ref{fig:phase2-structural} provides a useful explanation. Strong improvements are most often associated with introducing \texttt{MATERIALIZED} CTEs and reducing subquery count, both of which can simplify repeated computation. By contrast, regressions are more often associated with adding more CTE structure or extra decomposition without sufficient simplification. In other words, additional structure is beneficial only when it meaningfully improves the physical plan.

This interpretation is consistent with the omitted per-query analyses: the gains are concentrated in a limited subset of queries rather than spread evenly across the workload, and the distribution is strongly skewed. The key lesson is therefore not that the selected configuration is universally superior, but that it is effective when the model identifies a genuinely favorable structural rewrite.

\begin{figure}[htbp]
\centering
\includegraphics[width=0.72\linewidth]{datasets/artifacts/findings_figures/phase2_outcome_distribution.png}
\caption{Distribution of Phase~2 outcomes for the selected best configuration.}
\label{fig:phase2-outcomes}
\end{figure}

\begin{figure}[htbp]
\centering
\includegraphics[width=0.82\linewidth]{datasets/artifacts/findings_figures/phase2_structural_changes.png}
\caption{Most common structural rewrite patterns for strong improvements and regressions in Phase~2.}
\label{fig:phase2-structural}
\end{figure}

\subsection{Interpretation}

Overall, the empirical results show that the proposed pipeline can improve performance for a subset of TPC-DS queries, and that richer prompt context with direct generation is the strongest tested configuration. At the same time, the median Phase~2 improvement is small, several valid rewrites still regress, and some selected candidates remain invalid. These results reinforce the need for strict semantic validation and cautious performance-based selection when applying LLMs to SQL optimization.
